% TikZ settings
\tikzset{
  >=latex,
  every node/.style={thick, circle, draw, minimum height=1.7em, inner sep=0pt, text centered},
  every path/.style={thick},
}

\usetikzlibrary{calc, shapes.geometric, fadings, decorations.pathreplacing}

% Some custom colors
\definecolor{Crimson}{RGB}{165, 28, 48}
\definecolor{SteelBlue}{RGB}{70, 130, 180}
\definecolor{DarkRed}{RGB}{165, 0, 38}
\definecolor{LightRed}{RGB}{215, 48, 39}
\definecolor{DarkOrange}{RGB}{244, 109, 67}
\definecolor{LightOrange}{RGB}{253, 174, 97}
\definecolor{MidYellow}{RGB}{254, 224, 144}
\definecolor{MiddleBlue}{RGB}{116, 173, 209}
\definecolor{SkyBlue}{RGB}{69, 117, 180}
\definecolor{SeaBlue}{RGB}{49, 54, 149}
\definecolor{Pink}{HTML}{E91E63}
\definecolor{DeepOrange}{HTML}{FF5722}
\definecolor{Blue}{HTML}{2196F3}
\definecolor{Orange}{HTML}{FF9800}
\definecolor{Cyan}{HTML}{00BCD4}
\definecolor{Green}{HTML}{4CAF50}
\definecolor{DeepPurple}{HTML}{673AB7}
\definecolor{Purple}{HTML}{9C27B0}
\definecolor{Red}{HTML}{F44336}
\definecolor{Brown}{HTML}{795548}
\definecolor{Teal}{HTML}{009688}
\definecolor{Yellow}{HTML}{FFEB3B}
\definecolor{Grey}{HTML}{9E9E9E}
\definecolor{Lime}{HTML}{CDDC39}
\definecolor{Indigo}{HTML}{3F51B5}
\definecolor{LightGreen}{HTML}{8BC34A}
\definecolor{BlueGrey}{HTML}{607D8B}
\definecolor{LightBlue}{HTML}{03A9F4}
\definecolor{Amber}{HTML}{FFC107}

% special sets
\newcommand{\bN}{\mathbb{N}}
\newcommand{\bZ}{\mathbb{Z}}
\newcommand{\bQ}{\mathbb{Q}}
\newcommand{\bR}{\mathbb{R}}
\newcommand{\bC}{\mathbb{C}}

% special objects
\newcommand{\cR}{\mathcal{R}}
\newcommand{\cN}{\mathcal{N}}
\newcommand{\cG}{\mathcal{G}}
\newcommand{\tcG}{\tilde{\mathcal{G}}}
\newcommand{\bP}{\mathbb{P}}
\newcommand{\hY}{R}
\newcommand{\hy}{r}
\newcommand{\hR}{R}

% endings for equations
\def\eqp{\: .}
\def\eqc{\: ,}

% text surrounded by space
\newcommand{\qtxtq}[1]{\quad\text{#1}\quad}
\newcommand{\qqtxtqq}[1]{\qquad\text{#1}\qquad}

% shortcuts for frequently used terms
\newcommand*{\ie}{i.e.\@\xspace}
\newcommand*{\wrt}{w.r.t.\@\xspace}
\newcommand*{\etc}{etc.\@\xspace}

% references (just wrappers of usual \ref for flexibility)
\newcommand{\secref}[1]{Section~\ref{#1}}
\newcommand{\Secref}[1]{Section~\ref{#1}}
\newcommand{\figref}[1]{Figure~\ref{#1}}
\newcommand{\Figref}[1]{Figure~\ref{#1}}
\newcommand{\tabref}[1]{Table~\ref{#1}}
\newcommand{\Tabref}[1]{Table~\ref{#1}}

% shortcuts
\newcommand{\usub}[2]{\underset{#1}{\underbrace{#2}}}
\newcommand\given{{\,|\,}}
\newcommand\BB[1]{{\mathbb{#1}}}
\newcommand\B[1]{{\bm{#1}}}
\newcommand\C[1]{{\mathcal{#1}}}
\renewcommand{\dot}[2]{\langle #1, #2 \rangle}
\DeclareMathOperator*{\argmin}{arg\,min}
\DeclareMathOperator*{\argmax}{arg\,max}
\DeclareMathOperator{\E}{\mathbb{E}}
\DeclareMathOperator{\indep}{\perp\!\!\!\perp}
\DeclareMathOperator{\dep}{\not\! \perp\!\!\!\perp}

% math environments
\newtheorem{theorem}{Theorem}
\newtheorem*{theorem*}{Theorem}

\newtheorem{lemma}{Lemma}
\newtheorem*{lemma*}{Lemma}

\newtheorem{corollary}{Corollary}
\newtheorem*{corollary*}{Corollary}

\newtheorem{proposition}{Proposition}
\newtheorem*{proposition*}{Proposition}

\theoremstyle{definition}
\newtheorem{definition}{Definition}

% biblatex setup
\usepackage[%
  backend=bibtex,
  style=numeric-comp,
  hyperref=true,
  sorting=none,
  %eprint=true,
  maxbibnames=99,
  %doi=false,
  url=false,
  isbn=false,
  hyperref
  ]{biblatex}
\addbibresource{refs.bib}

